\protect \section *{\nameref *{sec:discharge}}
\begin{Solution}{11.{1}}
		$k_{iz}=C\bar v_e(\varepsilon_{iz}+2T_e)e^{-\varepsilon_{iz}/T_e}$; зростання приблизно в~$\approx6$ разів
	(за апроксимацією з~розд.~\ref{sec:ionization-rate}~--- у~$\approx6{,}5$ раза).
	
\end{Solution}
\begin{Solution}{11.{2}}
		$(pd)_{\min}=e\ln(1+1/\gamma)/A\approx0{,}84$~Торр$\,\cdot\,$см, $V_{\min}=B(pd)_{\min}\approx305$~В;
	при $pd=76$~Торр$\,\cdot\,$см $V_b\approx4{,}6$~кВ.
	
\end{Solution}
\begin{Solution}{11.{3}}
		$r_D\approx0{,}13$~мм, $s\approx3{,}2$~мм; кінцева швидкість іона $\approx3{,}8\times10^6$~см/с,
	час прольоту $3s/v\approx0{,}25$~мкс, тобто $\approx3{,}4$ періоду поля.
	
\end{Solution}
\begin{Solution}{11.{4}}
		$T_e\approx2{,}8$~еВ, $n_0\approx1{,}5\times10^{11}~\mathrm{см}^{-3}$; від потужності
	$T_e$ не залежить, а~зменшення тиску вдвічі підвищує її до $\approx3{,}3$~еВ; густина зросте менш ніж
	удвічі, бо зростають $u_B$ і~$\mathcal E_T$ ($A_{\rm eff}$ вважаємо сталою).
	
\end{Solution}
\begin{Solution}{11.{5}}
		$B=2\pi fmc/e\approx875$~Гс; $n_c\approx7{,}4\times10^{10}~\mathrm{см}^{-3}$; при $\omega<\omega_{ce}$
	права хвиля є свистовою і~не має відсічки за густиною.
	
\end{Solution}
\begin{Solution}{11.{6}}
		$k_\perp\approx1{,}5~\mathrm{см}^{-1}$, $k_z\approx0{,}31~\mathrm{см}^{-1}$, $k\approx1{,}6~\mathrm{см}^{-1}$;
	$n\approx1{,}4\times10^{13}~\mathrm{см}^{-3}$.
	
\end{Solution}
\begin{Solution}{11.{7}}
		$\rho_e\approx0{,}15$~мм (замагнічені), $\rho_i\approx7$~см (незамагнічені).
	
\end{Solution}
